\scanpage{35}

pour les $n'<n$. \deleted{Prouvons que} Nous savons donc déjà que \deleted{\unsure{$\ast$}}$(\ast^{\unsure{i}})$ est bijectif pour $i\leq n-1$, injectif pour $i=n$, on veut gagner un cran.

On a une suite exacte

\[
0\mathbin{\unsure{-|\to}}F\mathbin{-}L\mathbin{-}G\to0
\]

où \inserted{$G$ loc libre,} $L$ est une somme directe de modules loc. libres \deleted{\unsure{sommes|sur un}} $\simeq L_\alpha$, donc les assertions : pour $n$ sont vraies pour $L$.

Considérons l'hom de suites exactes

\diagram{\begin{tikzcd}[ampersand replacement=\&,column sep=small]
H^i(L) \arrow[r,dash] \arrow[d,"\wr\;\alpha"] \& H^i(G) \arrow[r,dash] \arrow[d,"\beta"] \& H^{i+1}(F) \arrow[r,dash] \arrow[d,"\varphi"] \& H^{i+1}(L) \arrow[r,dash] \arrow[d,"\wr\;\alpha'"] \& H^{i+1}(G) \arrow[d,"\beta'"] \\
H^i(\hat L) \arrow[r,dash] \& H^i(\hat G) \arrow[r,dash] \& H^{i+1}(\hat F) \arrow[r,dash] \& H^{i+1}(\hat L) \arrow[r,dash] \& H^{i+1}(\hat{\unsure{C|G}})
\end{tikzcd}}

pour $i=n-1,n$ donc $i+1=n,n+1$. Dans le premier cas, on veut prouver que la flèche médiane est surjective, dans le deuxième qu'elle est injective.

a) $i=n-1$, on sait par hyp. de réc. que $\beta$ est surjective, \deleted{(\unsure{est|ou|tout} bijective)} \deleted{il en est de même de $\alpha$, \unsure{enfin|alors} par l'hyp. de réc.} \inserted{de plus $\alpha'$ est surj par hyp. et $\beta'$ est injective.} Par lemme des 5, $\varphi$ est \deleted{bijective} surjective \inserted{(\deleted{\unsure{même|donc}} bijective)}

b) $\unsure{v|i}=n$, On sait par \deleted{a) l'hyp} \deleted{a) que $\beta$ est} \inserted{a)} \unsure{précemment|ci-dessus} que $\alpha$ et $\beta$ sont bijectives, ou $\alpha$ est surj, $\beta$ inj, de plus $\alpha'$ est inj.

% archive sheet number: 34
